% 第10章 银行CRM系统全栈开发
% 智慧银行实验教程——AI驱动的金融科技实践

\chapter{第10章\quad 银行CRM系统全栈开发}
\chapteroverview{
本章学习目标：
\begin{enumerate}
  \item 完成银行CRM系统的完整需求分析
  \item 设计并实现响应式前端界面
  \item 实现客户管理CRUD功能
  \item 构建基于规则的产品推荐模块
  \item 实现报表与数据可视化
  \item 完成系统测试与交付
\end{enumerate}
}

\keyterms{CRM系统、客户关系管理、Tailwind CSS、CRUD、localStorage、产品推荐、规则引擎、数据可视化、Chart.js、响应式设计、交叉销售}

% ============================================================
\section{银行CRM系统需求分析}
% ============================================================

\subsection{CRM系统概述}

客户关系管理（Customer Relationship Management，CRM）是银行业信息化的核心系统之一。在银行语境下，CRM系统的核心目标是帮助客户经理\textbf{更好地了解客户、服务客户、拓展客户}，从而提升客户满意度和银行收益。\par

据麦肯锡研究显示，有效使用CRM系统的银行，其交叉销售成功率可提升20\%--30\%，客户流失率降低15\%--25\%。在国内银行业数字化转型的大趋势下，CRM系统已从"可选工具"升级为"必备基础设施"。\par

银行CRM与通用CRM（如Salesforce、HubSpot）的关键区别在于：银行CRM需要深度整合客户资产数据（AUM）、风险评级、产品持有等专业信息，而不仅仅是联系人和销售管道管理。具体而言，银行CRM有以下专业特性：\par

\begin{itemize}[leftmargin=2em]
  \item \textbf{AUM导向}：以客户管理资产规模（Assets Under Management）为核心指标，而非销售额
  \item \textbf{风险合规}：所有产品推荐必须符合客户风险承受能力（适当性管理要求）
  \item \textbf{监管约束}：客户数据的采集、存储、使用必须符合《个人信息保护法》和银保监会相关规定
  \item \textbf{多产品交叉}：银行产品线覆盖存款、理财、基金、保险、贷款等，交叉销售空间大
\end{itemize}\par

一个优秀的银行CRM系统能够帮助客户经理实现以下目标：

\begin{itemize}[leftmargin=2em]
  \item \textbf{全面了解客户}：一目了然地查看客户的基本信息、资产状况、产品持有、交易历史
  \item \textbf{精准推荐产品}：基于客户画像，智能推荐适合的银行产品
  \item \textbf{高效管理客户}：快速搜索、筛选、分群，提高客户管理效率
  \item \textbf{数据驱动决策}：通过数据看板了解整体客户结构，发现业务机会
\end{itemize}

\subsection{目标用户画像}

理解目标用户是做好需求分析的前提。在真实项目中，产品团队会通过用户访谈、问卷调查、现场观察等方法构建详细的用户画像。在我们的教学实验中，我们将通过预设的用户画像来模拟这一过程。

\begin{table}[H]
\centering
\begin{tabular}{lp{9cm}}
\toprule
\textbf{维度} & \textbf{描述} \\
\midrule
角色 & 支行/网点一线客户经理 \\
年龄 & 25--40岁 \\
技术水平 & 日常使用办公软件，非技术背景 \\
核心痛点 & 客户信息分散在多个系统、产品推荐靠经验、缺乏数据支撑 \\
使用场景 & 晨会准备、客户拜访前查询、面谈中记录、周报汇总 \\
时间压力 & 单日需管理50--200名客户，每名客户平均接触时间有限 \\
期望功能 & 一键查看客户全景、智能推荐产品、自动生成分析报告 \\
\bottomrule
\end{tabular}
\caption{银行客户经理用户画像}
\end{table}

\begin{casebox}[典型使用场景]
\textbf{场景}：客户经理李明早上8:30到支行，9:00约了一位金卡客户王女士来办理业务。\par

在王女士到达之前，李明打开CRM系统，搜索王女士的名字，立即看到她的全景信息：AUM 280万元，持有3款产品（活期宝、稳盈债券、平衡配置），风险等级为平衡型，上次联系是两周前。系统推荐了2款新产品：安心固收（稳健收益，与已持有债券基金互补）和优选成长（收益更高，风险仍在可承受范围内）。李明快速浏览推荐理由和话术，为9:00的会面做好准备。\par

\textbf{如果没有CRM}，李明需要分别登录核心系统查账户、理财系统查产品、Excel查客户资料，耗时且容易遗漏信息。\end{casebox}

基于此用户画像，我们的CRM系统设计原则是：\textbf{简单直观、一键操作、信息聚合}——客户经理不需要培训就能上手，关键操作不超过3次点击。

\subsection{核心功能模块划分}

基于用户画像和业务需求分析，银行CRM系统划分为以下五大功能模块：

\begin{table}[H]
\centering
\small
\begin{tabular}{llp{6cm}}
\toprule
\textbf{模块} & \textbf{优先级} & \textbf{核心功能} \\
\midrule
客户管理 & Must & 查看/搜索/新增/编辑/分群标签 \\
产品管理 & Should & 产品库浏览/客户持有产品查看 \\
智能推荐 & Must & 基于画像的产品推荐+推荐理由+话术 \\
报表分析 & Should & 客户分布/产品渗透率/KPI看板 \\
工作台 & Could & 待办事项/日历/消息提醒 \\
\bottomrule
\end{tabular}
\caption{CRM功能模块与优先级}
\end{table}

\subsection{用户故事列表}

基于功能模块，我们编写了12条用户故事：

\begin{table}[H]
\centering
\small
\begin{tabular}{clp{7.5cm}}
\toprule
\textbf{编号} & \textbf{模块} & \textbf{用户故事} \\
\midrule
US-001 & 客户管理 & As a 客户经理, I want 查看客户列表, so that 我能快速了解所有客户概况 \\
US-002 & 客户管理 & As a 客户经理, I want 按姓名/VIP/手机号搜索客户, so that 我能快速定位目标客户 \\
US-003 & 客户管理 & As a 客户经理, I want 查看客户详细信息及持有产品, so that 我能全面了解客户资产状况 \\
US-004 & 客户管理 & As a 客户经理, I want 新增客户信息, so that 我能在首次面谈后快速录入 \\
US-005 & 客户管理 & As a 客户经理, I want 编辑客户信息, so that 我能及时更新客户资料 \\
US-006 & 产品推荐 & As a 客户经理, I want 查看客户的产品推荐, so that 我能精准推荐适合的产品 \\
US-007 & 产品推荐 & As a 客户经理, I want 看到推荐理由和话术, so that 我能更有说服力地向客户介绍 \\
US-008 & 产品管理 & As a 客户经理, I want 浏览银行产品库, so that 我能了解所有可售产品 \\
US-009 & 报表分析 & As a 客户经理, I want 查看客户AUM分布图, so that 我能了解整体客户资产结构 \\
US-010 & 报表分析 & As a 客户经理, I want 查看VIP等级占比, so that 我能评估客户分层效果 \\
US-011 & 报表分析 & As a 客户经理, I want 查看产品渗透率, so that 我能发现交叉销售机会 \\
US-012 & 工作台 & As a 客户经理, I want 看到待办事项和消息提醒, so that 我不会遗漏重要工作 \\
\bottomrule
\end{tabular}
\caption{银行CRM用户故事列表}
\end{table}

\subsection{功能优先级排序}

受课堂时间限制，本实验实现Top 6核心功能：

\begin{table}[H]
\centering
\begin{tabular}{cll}
\toprule
\textbf{实现顺序} & \textbf{用户故事} & \textbf{理由} \\
\midrule
1 & US-001 客户列表 & 最基础的功能，其他功能都依赖它 \\
2 & US-003 客户详情 & 核心查看功能 \\
3 & US-004 新增客户 & 核心录入功能 \\
4 & US-002 搜索筛选 & 提升效率的关键功能 \\
5 & US-006 产品推荐 & 差异化价值最大的功能 \\
6 & US-009 数据看板 & 数据可视化展示 \\
\bottomrule
\end{tabular}
\caption{本实验实现的Top 6功能}
\end{table}

\begin{notebox}[关于功能范围控制]
很多学生在开发过程中容易"功能蔓延"——实现了一个功能后忍不住再加一个，结果核心功能反而没做好。请严格遵守上述Top 6的优先级顺序：只有当前功能完全实现并通过测试后，才进入下一个功能。如果还有余力，优先实现US-005（编辑客户）和US-010（VIP占比图），而非跳跃到更花哨的功能。\end{notebox}

\begin{theorybox}[需求分析中的常见陷阱]
在BMAD的B阶段，学生常犯以下错误：
\begin{itemize}[leftmargin=2em]
  \item \textbf{功能堆砌}：列出太多功能，缺乏优先级排序。结果是每个功能都只做到一半，没有一个完整的
  \item \textbf{忽视非功能需求}：只关注"做什么"，忽略了"做得怎么样"——性能、可用性、容错性
  \item \textbf{用户故事缺少价值}：写成"As a 客户经理, I want 搜索功能"，缺少"so that"部分——没有说明功能的价值，后续难以判断优先级
  \item \textbf{验收标准模糊}：写"搜索功能正常"而非"输入姓名关键字后1秒内返回匹配结果"
\end{itemize}
好的需求分析是\textbf{具体、可验证、有优先级}的。在BMAD中，B阶段的质量直接决定了后续所有阶段的效率。\end{theorybox}

% ============================================================
\section{前端界面设计}
% ============================================================

\subsection{Tailwind CSS快速入门}

Tailwind CSS是一个实用类优先（utility-first）的CSS框架，与传统的组件级框架（如Bootstrap）不同，Tailwind提供的是\textbf{原子化的CSS类}——每个类只做一件事，通过组合这些类来构建界面。

\textbf{实用类设计哲学}：传统CSS是"先写HTML，再写CSS样式"；Tailwind是"在HTML中直接写样式类"。这种方式的优点是\textbf{不需要在CSS文件和HTML文件之间来回切换}，开发速度更快，且不会产生冗余的CSS代码。

\textbf{常用类速查}：

\begin{table}[H]
\centering
\small
\begin{tabular}{lll}
\toprule
\textbf{类别} & \textbf{常用类} & \textbf{效果} \\
\midrule
布局 & \texttt{flex}, \texttt{grid}, \texttt{container} & 弹性/网格/容器布局 \\
间距 & \texttt{p-4}, \texttt{m-2}, \texttt{gap-3} & 内边距/外边距/间距 \\
排版 & \texttt{text-xl}, \texttt{font-bold}, \texttt{text-center} & 字号/字重/对齐 \\
颜色 & \texttt{bg-blue-500}, \texttt{text-gray-700} & 背景色/文字色 \\
边框 & \texttt{border}, \texttt{rounded-lg}, \texttt{shadow-md} & 边框/圆角/阴影 \\
响应式 & \texttt{md:flex}, \texttt{lg:grid-cols-3} & 中屏/大屏断点 \\
\bottomrule
\end{tabular}
\caption{Tailwind CSS常用类速查}
\end{table}

\begin{notebox}[Tailwind CSS的CDN引入方式]
本课程使用CDN方式引入Tailwind CSS，无需安装任何依赖。只需在HTML的\texttt{<head>}中添加一行：

\begin{lstlisting}[style=html]
<script src="https://cdn.tailwindcss.com"></script>
\end{lstlisting}

这种方式适合原型开发和教学演示。在生产环境中，推荐使用npm安装并配置构建流程，以获得更小的文件体积和更好的性能。
\end{notebox}

\textbf{银行蓝色系配色方案}：金融行业的UI设计有其独特的配色传统——以蓝色为主色调，传达专业、稳健、可信赖的品牌形象。本系统的配色方案如下：
\begin{itemize}[leftmargin=2em]
  \item \textbf{主色}：\texttt{\#1e40af}（深蓝色）——导航栏背景、重要按钮
  \item \textbf{辅色}：\texttt{\#3b82f6}（中蓝色）——侧边栏选中态、链接、强调色
  \item \textbf{浅色}：\texttt{\#eff6ff}（极浅蓝）——页面背景、卡片背景
  \item \textbf{成功}：\texttt{\#10b981}（绿色）——操作成功提示、增长指标
  \item \textbf{警告}：\texttt{\#f59e0b}（橙色）——风险提示、注意事项
  \item \textbf{错误}：\texttt{\#ef4444}（红色）——表单验证失败、危险操作
\end{itemize}

在Tailwind中，这些颜色可以通过\texttt{tailwind.config}自定义，或直接使用Tailwind内置的蓝色系（\texttt{blue-50}到\texttt{blue-900}）。本课程使用内置色值即可满足需求。

\subsection{页面结构设计}

银行CRM采用经典的\textbf{侧边栏布局}，页面从上到下、从左到右分为四个区域：

\textbf{导航栏（顶部，高60px）}：
\begin{itemize}[leftmargin=2em]
  \item 左侧：Logo图标 + "智慧银行CRM"系统名称
  \item 右侧：当前用户姓名 + 退出按钮
  \item 样式：深蓝色背景（bg-blue-800）+ 白色文字
\end{itemize}

\textbf{侧边栏（左侧，宽200px）}：
\begin{itemize}[leftmargin=2em]
  \item 功能菜单列表，每项带图标：
  \item 客户管理（默认选中）
  \item 产品推荐
  \item 报表分析
  \item 工作台
  \item 样式：浅蓝色背景（bg-blue-50）+ 选中项高亮
\end{itemize}

\textbf{主内容区（中间，自适应宽度）}：
\begin{itemize}[leftmargin=2em]
  \item 根据侧边栏选择动态切换：
  \item 客户管理→客户列表/客户详情/新增客户表单
  \item 产品推荐→推荐列表
  \item 报表分析→数据看板
  \item 工作台→待办列表
\end{itemize}

\textbf{底部栏（底部，高40px）}：
\begin{itemize}[leftmargin=2em]
  \item 居中显示版权信息："© 2026 智慧银行CRM系统 | 四川农业大学经济学院"
  \item 样式：灰色背景 + 小字号
\end{itemize}

\subsection{响应式设计策略}

本CRM系统需要兼容桌面和平板两种屏幕尺寸：

\begin{table}[H]
\centering
\begin{tabular}{lll}
\toprule
\textbf{屏幕} & \textbf{断点} & \textbf{布局调整} \\
\midrule
桌面 & $\geq$ 1024px & 侧边栏常驻显示，主内容区自适应宽度 \\
平板 & 768--1023px & 侧边栏折叠为汉堡菜单，主内容区全宽 \\
手机 & $<$ 768px & 底部导航栏替代侧边栏，列表单列显示 \\
\bottomrule
\end{tabular}
\caption{响应式断点设计}
\end{table}

在Tailwind中，响应式通过前缀实现：\texttt{md:}表示768px以上生效，\texttt{lg:}表示1024px以上生效。例如：

\begin{lstlisting}[style=html]
<!-- 桌面端侧边栏常驻，平板端隐藏 -->
<aside class="hidden lg:block w-48 bg-blue-50">
  <!-- 侧边栏内容 -->
</aside>
<!-- 平板端汉堡菜单按钮 -->
<button class="lg:hidden" onclick="toggleSidebar()">
  ☰
</button>
\end{lstlisting}

\begin{practicebox}[用AI设计CRM界面]
\textbf{任务}：使用AI生成银行CRM的HTML骨架页面。

\textbf{提示词}：

\begin{lstlisting}
请基于 design.md 生成银行CRM的HTML骨架页面，要求：
1. 单文件 web/index.html
2. 引入 Tailwind CSS CDN
3. 页面结构：导航栏(深蓝) + 侧边栏(浅蓝) + 主内容区 + 底部栏
4. 配色方案：银行蓝色系（主色 #1e40af，辅色 #3b82f6，背景 #eff6ff）
5. 侧边栏菜单项：客户管理/产品推荐/报表分析/工作台
6. 主内容区默认显示"欢迎使用智慧银行CRM系统"
7. 响应式：桌面端侧边栏常驻，平板端侧边栏可折叠
8. 不需要实现具体功能逻辑，只需页面骨架和样式
\end{lstlisting}

\textbf{审核要点}：页面在浏览器中打开后布局是否正确？配色是否统一？侧边栏菜单项是否完整？响应式切换是否正常？
\end{practicebox}

% ============================================================
\section{客户管理模块}
% ============================================================

\subsection{数据模型设计}

客户对象是CRM系统的核心数据实体，其字段设计需要兼顾业务需求和技术实现：

\begin{table}[H]
\centering
\small
\begin{tabular}{llllp{4cm}}
\toprule
\textbf{分组} & \textbf{字段名} & \textbf{类型} & \textbf{必填} & \textbf{说明} \\
\midrule
\multirow{5}{*}{基本信息} & customerId & string & 自动 & UUID自动生成 \\
 & name & string & 是 & 客户姓名 \\
 & phone & string & 是 & 手机号（11位） \\
 & idCard & string & 否 & 身份证号（18位） \\
 & address & string & 否 & 通讯地址 \\
\midrule
\multirow{4}{*}{银行信息} & branch & string & 是 & 开户行 \\
 & accountNo & string & 自动 & 银行账号自动生成 \\
 & vipLevel & string & 是 & VIP等级：普通/银卡/金卡/白金/钻石 \\
 & managerName & string & 是 & 客户经理姓名 \\
\midrule
\multirow{4}{*}{资产信息} & aum & number & 是 & 管理资产规模（万元） \\
 & riskLevel & string & 是 & 风险等级：保守/稳健/平衡/进取/激进 \\
 & holdings & array & 自动 & 持有产品列表 \\
 & createdAt & string & 自动 & 创建时间 \\
\bottomrule
\end{tabular}
\caption{客户对象字段定义}
\end{table}

\subsection{CRUD功能实现}

CRUD（Create/Read/Update/Delete）是数据操作的四个基本动作，是所有管理系统的核心功能。在前端应用中实现CRUD，核心是设计好\textbf{数据操作函数}和\textbf{页面渲染函数}的分离——数据函数只负责读写localStorage，渲染函数只负责将数据转化为HTML，两者通过数据流转衔接。

\subsubsection{Create：新增客户}

新增客户表单包含以下字段和验证规则：

\begin{itemize}[leftmargin=2em]
  \item 姓名：必填，长度2--20
  \item 手机号：必填，正则验证\texttt{\^{}1[3-9]\textbackslash d\{9\}\$}
  \item 身份证号：选填，18位格式验证
  \item 开户行：必填，下拉选择
  \item VIP等级：必填，单选按钮组
  \item AUM：必填，正数验证
  \item 风险等级：必填，下拉选择
  \item 客户经理：必填，当前登录用户
\end{itemize}

提交后的处理流程：
\begin{enumerate}[leftmargin=2em]
  \item 前端表单验证通过
  \item 自动生成customerId和accountNo
  \item 记录createdAt时间戳
  \item 初始化holdings为空数组
  \item 写入localStorage
  \item 刷新客户列表
  \item 显示成功提示
\end{enumerate}

\subsubsection{Read：客户列表与详情}

客户列表页面包含以下功能：

\textbf{列表展示}：
\begin{itemize}[leftmargin=2em]
  \item 默认按AUM降序排列
  \item 每行显示：姓名、VIP等级（带颜色标签）、AUM、手机号、风险等级
  \item 每页10条，底部显示分页控件
  \item 顶部显示客户总数和AUM合计
\end{itemize}

\textbf{搜索筛选}：
\begin{itemize}[leftmargin=2em]
  \item 关键词搜索：模糊匹配姓名和手机号
  \item VIP等级筛选：下拉选择（全部/普通/银卡/金卡/白金/钻石）
  \item 风险等级筛选：下拉选择
  \item 搜索与筛选可叠加使用
\end{itemize}

\textbf{客户详情}：
\begin{itemize}[leftmargin=2em]
  \item 点击列表行进入详情页
  \item 分区域展示：基本信息/银行信息/资产信息/持有产品/推荐产品
  \item 右上角"编辑"和"返回"按钮
\end{itemize}

\subsubsection{Update：编辑客户信息}

编辑功能与新增共享表单组件，区别在于：
\begin{itemize}[leftmargin=2em]
  \item 表单预填充当前客户数据
  \item customerId和createdAt不可修改
  \item 保存时覆盖原记录（基于customerId匹配）
  \item 操作后显示"修改成功"提示
\end{itemize}

\subsubsection{Delete：标记删除}

本系统采用\textbf{逻辑删除}（软删除）策略，而非物理删除：
\begin{itemize}[leftmargin=2em]
  \item 客户对象增加\texttt{isDeleted}字段，默认为\texttt{false}
  \item 删除操作将\texttt{isDeleted}设为\texttt{true}，而非从数组中移除
  \item 列表查询时自动过滤\texttt{isDeleted=true}的记录
  \item 逻辑删除的好处：数据可恢复、保留审计记录
\end{itemize}

\subsection{localStorage持久化方案}

localStorage是浏览器提供的本地存储API，数据以键值对形式保存在浏览器中，刷新页面不会丢失。本系统的持久化方案如下：

\begin{lstlisting}[style=javascript]
// 数据初始化：首次访问时加载模拟数据
function initData() {
  if (!localStorage.getItem('crm_customers')) {
    localStorage.setItem('crm_customers',
      JSON.stringify(mockCustomers));
  }
}

// 读取数据
function getCustomers() {
  return JSON.parse(
    localStorage.getItem('crm_customers') || '[]'
  ).filter(c => !c.isDeleted);
}

// 保存数据
function saveCustomers(customers) {
  localStorage.setItem('crm_customers',
    JSON.stringify(customers));
}

// 新增客户
function addCustomer(customer) {
  const customers = getCustomers();
  customer.customerId = generateUUID();
  customer.accountNo = '6217' +
    Math.floor(Math.random()*1e12);
  customer.createdAt = new Date().toISOString();
  customer.holdings = [];
  customer.isDeleted = false;
  customers.push(customer);
  saveCustomers(customers);
}
\end{lstlisting}

\begin{warningbox}[localStorage的容量限制]
localStorage的存储容量通常为5--10MB，足以存储数百条客户记录。但如果数据量超过限制，\texttt{setItem}会抛出\texttt{QuotaExceededError}异常。在生产系统中应使用try-catch处理此异常，并提示用户清理数据。本实验中5--10条模拟数据不会触发此限制。
\end{warningbox}

\begin{practicebox}[分步骤用AI实现客户管理]
\textbf{提示词1：客户列表页（含搜索和筛选）}

\begin{lstlisting}
请在 web/index.html 的主内容区实现客户列表页，要求：
1. 显示客户列表，字段：姓名/VIP等级(彩色标签)/AUM/手机号/风险等级
2. 默认按AUM降序排列
3. 顶部搜索栏：关键词搜索(模糊匹配姓名+手机号)
4. 顶部筛选栏：VIP等级下拉 + 风险等级下拉
5. 搜索和筛选可叠加使用
6. 每页10条，底部分页控件
7. 列表顶部显示客户总数和AUM合计
8. 点击某行跳转到客户详情页
9. 右上角"新增客户"按钮
10. 数据从localStorage读取，键名 crm_customers
11. 如果localStorage无数据，初始化>=5条模拟客户数据
12. 使用Tailwind CSS美化界面
\end{lstlisting}

\textbf{提示词2：新增客户表单}

\begin{lstlisting}
请在客户列表页的"新增客户"按钮点击后，在主内容区显示新增客户表单，要求：
1. 表单字段：姓名(必填)/手机号(必填,11位)/身份证号(选填,18位)
   /开户行(必填,下拉)/VIP等级(必填,单选)/AUM(必填,正数)
   /风险等级(必填,下拉)/客户经理(必填)/地址(选填)
2. 前端表单验证：手机号正则/^1[3-9]\d{9}$/，身份证正则
   /^\d{17}[\dXx]$/
3. 提交后自动生成customerId(UUID)和accountNo
4. 写入localStorage，刷新列表，显示成功提示
5. 表单样式：卡片布局，字段分组显示，Tailwind美化
6. 底部"保存"和"取消"按钮
\end{lstlisting}

\textbf{提示词3：客户详情页}

\begin{lstlisting}
请在客户列表页点击某行后，在主内容区显示客户详情页，要求：
1. 分区域展示客户信息：
   - 基本信息区：姓名/手机号/身份证号/地址
   - 银行信息区：开户行/账号/VIP等级/客户经理
   - 资产信息区：AUM/风险等级
   - 持有产品区：产品列表（产品名/金额/购买日期）
2. 右上角"编辑"和"返回列表"按钮
3. 编辑按钮跳转到编辑表单（预填当前数据）
4. 样式：卡片式布局，区域间有分隔，Tailwind美化
5. 数据从localStorage读取
\end{lstlisting}

\textbf{每个提示词的审核要点}：功能是否完整？表单验证是否生效？localStorage读写是否正确？页面样式是否美观？交互是否流畅？
\end{practicebox}

% ============================================================
\section{产品推荐模块}
% ============================================================

\subsection{推荐策略设计}

产品推荐是银行CRM的核心差异化功能。本系统采用\textbf{基于规则}的推荐策略，不需要机器学习模型，通过预定义的业务规则匹配客户画像与产品特征。

\textbf{规则1：基于风险偏好匹配}。每款银行产品有对应的风险等级（R1--R5），每个客户也有风险等级（保守--激进）。推荐策略是\textbf{产品的风险等级不超过客户的风险等级}。例如，保守型客户只推荐R1（低风险）产品，而进取型客户可以推荐R1--R4产品。

\textbf{规则2：基于AUM等级匹配}。部分产品有最低购买金额要求，推荐策略是\textbf{客户的AUM不低于产品的最低购买金额}。例如，某私募基金最低购买金额为100万元，只有AUM≥100万的客户才会被推荐此产品。

\textbf{规则3：基于已持有产品的交叉销售}。如果客户已持有某类产品（如货币基金），推荐同类别或互补类别的产品（如短期理财、债券基金）。交叉销售的逻辑是\textbf{客户已接受某类产品，更容易接受类似产品}。

三条规则的优先级：规则1（风险匹配）> 规则2（AUM匹配）> 规则3（交叉销售）。即先过滤风险不匹配的产品，再过滤买不起的产品，最后按交叉销售排序。

\subsection{产品库数据设计}

产品库是推荐引擎的基础数据，本系统设计了12款银行产品，覆盖主流产品类别：

\begin{table}[H]
\centering
\small
\begin{tabular}{llcccl}
\toprule
\textbf{产品名称} & \textbf{类别} & \textbf{风险等级} & \textbf{最低金额(万)} & \textbf{年化收益率} & \textbf{适合客户} \\
\midrule
活期宝 & 货币基金 & R1 & 0.01 & 2.1\% & 保守型 \\
天天理财 & 短期理财 & R1 & 1 & 2.8\% & 保守型 \\
稳盈债券 & 债券基金 & R2 & 5 & 3.5\% & 稳健型 \\
安心固收 & 固定收益 & R2 & 10 & 4.0\% & 稳健型 \\
平衡配置 & 混合基金 & R3 & 5 & 5.2\% & 平衡型 \\
优选成长 & 股票基金 & R3 & 10 & 6.8\% & 平衡型 \\
价值精选 & 股票基金 & R4 & 20 & 8.5\% & 进取型 \\
全球配置 & QDII基金 & R4 & 30 & 7.2\% & 进取型 \\
量化对冲 & 对冲基金 & R4 & 50 & 6.5\% & 进取型 \\
科创板主题 & 股票基金 & R5 & 50 & 12.0\% & 激进型 \\
新兴市场 & QDII基金 & R5 & 100 & 10.5\% & 激进型 \\
私募股权 & 私募基金 & R5 & 100 & 15.0\% & 激进型 \\
\bottomrule
\end{tabular}
\caption{银行产品库数据}
\end{table}

\subsection{推荐结果展示}

推荐结果以卡片形式展示，每张卡片包含：

\begin{itemize}[leftmargin=2em]
  \item \textbf{产品名称和类别}：醒目的标题
  \item \textbf{关键指标}：风险等级标签/年化收益率/最低金额
  \item \textbf{推荐理由}：匹配了哪些规则（如"风险偏好匹配""AUM满足最低要求""与已持有产品互补"）
  \item \textbf{客户经理话术}：针对该客户的一句推荐语（如"张先生，您目前持有的货币基金收益较低，这款债券基金风险相近但收益更高，非常适合您"）
\end{itemize}

推荐结果只展示\textbf{Top 3}产品，避免信息过载。排序依据：规则匹配数多的优先，匹配数相同则按年化收益率降序。\par

\begin{theorybox}[从规则引擎到智能推荐]
本系统采用基于规则的产品推荐，优点是逻辑透明、可解释性强、实现简单。但在实际银行系统中，推荐引擎可能采用更复杂的方法：
\begin{itemize}[leftmargin=2em]
  \item \textbf{协同过滤}：基于"相似客户也购买了"的逻辑推荐，类似电商的"猜你喜欢"
  \item \textbf{内容推荐}：基于产品特征与客户偏好的匹配度推荐
  \item \textbf{混合推荐}：结合规则引擎和机器学习模型，兼顾合规性和精准度
\end{itemize}
银行业的推荐系统有一个独特约束——\textbf{适当性管理}。根据银保监会规定，金融机构向客户推荐产品时，必须确保产品的风险等级不超过客户的风险承受能力。这意味着，无论推荐算法多么先进，\textbf{风险匹配规则是不可逾越的底线}。我们的规则1正是对这一监管要求的实现。\end{theorybox}

\begin{practicebox}[用AI实现推荐引擎]
\textbf{任务}：使用AI实现基于规则的产品推荐模块。

\textbf{提示词}：

\begin{lstlisting}
请在 web/index.html 中实现银行CRM的产品推荐功能，要求：

1. 在侧边栏点击"产品推荐"后，主内容区显示推荐页面
2. 页面顶部有客户选择下拉框，选择客户后自动生成推荐
3. 推荐算法（基于规则）：
   规则1-风险匹配：产品风险等级 <= 客户风险等级
     （保守=R1, 稳健=R1-R2, 平衡=R1-R3,
      进取=R1-R4, 激进=R1-R5）
   规则2-AUM匹配：客户AUM >= 产品最低金额
   规则3-交叉销售：客户已持有某类产品，推荐同类或互补类
4. 推荐结果只展示Top 3产品，每张产品卡片包含：
   - 产品名称/类别/风险等级(彩色标签)/年化收益率/最低金额
   - 推荐理由（匹配了哪些规则）
   - 客户经理话术（一句话推荐语）
5. 如果客户无推荐产品，显示"暂无适合该客户的产品"
6. 产品库数据定义在JS中，包含12款银行产品
7. 数据从localStorage读取客户信息和持有产品
8. 使用Tailwind CSS美化卡片样式
\end{lstlisting}

\textbf{审核要点}：推荐算法是否正确实现了三条规则？风险匹配是否准确？AUM过滤是否生效？Top 3排序是否合理？推荐理由和话术是否有针对性？
\end{practicebox}

% ============================================================
\section{报表与数据可视化模块}
% ============================================================

\subsection{图表库选择与引入}

本系统选择\textbf{Chart.js}作为图表库，理由如下：
\begin{itemize}[leftmargin=2em]
  \item 轻量级：压缩后仅60KB，不拖慢页面加载
  \item API简洁：几行代码即可创建图表
  \item 样式美观：默认配色和动画效果专业
  \item 响应式：图表自动适应容器尺寸
\end{itemize}

CDN引入方式：

\begin{lstlisting}[style=html]
<script src="https://cdn.jsdelivr.net/npm/chart.js@4"></script>
\end{lstlisting}

\begin{notebox}[Chart.js vs ECharts]
ECharts是百度开源的图表库，功能更强大（支持3D图表、地理图表、大数据量渲染），但体积更大（1MB+）且API更复杂。本课程选择Chart.js是因为它的学习曲线更平缓，适合快速原型开发。如果你未来需要更丰富的图表类型，可以平滑迁移到ECharts。
\end{notebox}

\subsection{关键报表设计}

本系统设计了4个核心数据看板：

\textbf{报表1：客户AUM分布饼图}。将客户按AUM区间分为5档：0--50万、50--200万、200--500万、500--1000万、1000万以上。饼图直观展示各档的客户数量占比，帮助客户经理了解客户资产结构。

\textbf{报表2：月度新增客户柱状图}。按月统计新增客户数量，柱状图展示近6个月的趋势变化。如果数据不足6个月，可以用模拟数据填充历史月份。

\textbf{报表3：VIP等级占比环形图}。环形图展示各VIP等级（普通/银卡/金卡/白金/钻石）的客户数量占比，核心指标放在环形中央（如客户总数）。

\textbf{报表4：产品渗透率热力图}。横轴为VIP等级，纵轴为产品类别，颜色深浅表示该等级客户持有该类别产品的比例。热力图帮助发现交叉销售机会——颜色浅的格子表示渗透率低，存在提升空间。

\subsection{数据聚合逻辑}

图表数据需要从客户原始数据中聚合计算。以下是核心的聚合函数：

\begin{lstlisting}[style=javascript]
// AUM分布聚合
function getAUMDistribution(customers) {
  const ranges = [
    {label:'0-50万', min:0, max:50},
    {label:'50-200万', min:50, max:200},
    {label:'200-500万', min:200, max:500},
    {label:'500-1000万', min:500, max:1000},
    {label:'1000万+', min:1000, max:Infinity}
  ];
  return ranges.map(r => ({
    label: r.label,
    count: customers.filter(c =>
      c.aum >= r.min && c.aum < r.max).length
  }));
}

// VIP等级聚合
function getVIPDistribution(customers) {
  const levels = ['普通','银卡','金卡','白金','钻石'];
  return levels.map(l => ({
    label: l,
    count: customers.filter(c =>
      c.vipLevel === l).length
  }));
}
\end{lstlisting}

\begin{practicebox}[用AI实现数据看板]
\textbf{任务}：使用AI实现银行CRM的报表与数据可视化模块。

\textbf{提示词}：

\begin{lstlisting}
请在 web/index.html 中实现银行CRM的报表分析页面，要求：

1. 侧边栏点击"报表分析"后，主内容区显示数据看板
2. 看板包含4个图表，2行2列网格布局：
   - 左上：客户AUM分布饼图
     分5档：0-50万/50-200万/200-500万/500-1000万/1000万+
   - 右上：VIP等级占比环形图
     5个等级：普通/银卡/金卡/白金/钻石
   - 左下：月度新增客户柱状图
     近6个月趋势，数据从客户createdAt字段统计
   - 右下：产品持有类别分布柱状图
     按产品类别统计持有客户数
3. 每个图表上方显示标题和简要说明
4. 图表使用Chart.js渲染，配色与系统蓝色主题一致
5. 数据从localStorage的crm_customers中聚合计算
6. 如果数据不足，显示模拟数据作为示例
7. 顶部显示3个KPI卡片：客户总数/AUM总计/平均AUM
8. 使用Tailwind CSS美化布局
9. 图表容器有白色背景和阴影，增加视觉层次
\end{lstlisting}

\textbf{审核要点}：图表是否正确渲染？数据聚合逻辑是否正确？配色是否统一？KPI卡片数据是否准确？图表在窗口缩放时是否自适应？
\end{practicebox}

% ============================================================
\section{系统测试与部署}
% ============================================================

\subsection{功能测试清单}

系统开发完成后，需要逐项验证每个模块的功能。测试不是可选项，而是软件交付的必要环节——未经测试的代码不能称为"完成"。以下是按模块组织的测试清单：

\begin{table}[H]
\centering
\small
\begin{tabular}{clp{5.5cm}l}
\toprule
\textbf{模块} & \textbf{编号} & \textbf{测试项} & \textbf{预期结果} \\
\midrule
\multirow{4}{*}{客户管理} & T-001 & 打开页面显示客户列表 & 列表正常渲染，数据完整 \\
 & T-002 & 按姓名搜索客户 & 返回匹配结果 \\
 & T-003 & 新增客户（完整信息） & 客户成功添加 \\
 & T-004 & 新增客户（姓名为空） & 提示"请输入姓名" \\
 & T-005 & 点击客户查看详情 & 显示完整信息 \\
\midrule
\multirow{3}{*}{产品推荐} & T-006 & 选择客户生成推荐 & 显示Top 3产品 \\
 & T-007 & 保守型客户推荐 & 只显示R1产品 \\
 & T-008 & 无推荐结果提示 & 显示"暂无适合产品" \\
\midrule
\multirow{3}{*}{报表分析} & T-009 & AUM分布饼图渲染 & 图表正常显示 \\
 & T-010 & KPI卡片数据 & 数值与列表数据一致 \\
 & T-011 & 图表响应式 & 窗口缩放时图表自适应 \\
\midrule
\multirow{2}{*}{持久化} & T-012 & 刷新页面数据保留 & localStorage数据不丢失 \\
 & T-013 & 新增后刷新验证 & 新增客户在刷新后仍存在 \\
\bottomrule
\end{tabular}
\caption{银行CRM功能测试清单}
\end{table}

\subsection{常见bug与AI修复循环}

在测试过程中，你可能会遇到以下常见bug。对于每个bug，我们都提供了标准的描述格式，方便你直接发给AI修复：

\textbf{Bug 1：客户列表不显示}
\begin{itemize}[leftmargin=2em]
  \item 复现步骤：打开页面，主内容区空白
  \item 期望行为：显示客户列表
  \item 可能原因：localStorage初始化失败，或渲染函数未调用
  \item 修复提示词：
\end{itemize}

\begin{lstlisting}
我的CRM系统打开页面后客户列表为空白。
请检查以下可能的问题并修复：
1. localStorage是否正确初始化了模拟数据
2. 页面加载时是否调用了渲染函数
3. 数据读取函数是否有错误
4. HTML中是否有正确的容器元素
请提供修复后的完整代码。
\end{lstlisting}

\textbf{Bug 2：新增客户后列表不刷新}
\begin{itemize}[leftmargin=2em]
  \item 复现步骤：填写新增客户表单，点击保存，列表未更新
  \item 期望行为：列表立即显示新增的客户
  \item 可能原因：保存后未重新渲染列表，或渲染函数读取了缓存数据
\end{itemize}

\textbf{Bug 3：搜索功能无反应}
\begin{itemize}[leftmargin=2em]
  \item 复现步骤：在搜索框输入文字，列表未变化
  \item 期望行为：列表实时过滤匹配结果
  \item 可能原因：搜索框的事件监听器未绑定，或过滤函数逻辑有误
\end{itemize}

\textbf{Bug 4：刷新后数据丢失}
\begin{itemize}[leftmargin=2em]
  \item 复现步骤：新增客户后刷新页面，新增的客户消失了
  \item 期望行为：刷新后数据与刷新前一致
  \item 可能原因：新增时未写入localStorage，或初始化函数覆盖了已有数据
\end{itemize}

\subsection{Live Server本地预览}

Live Server是VS Code的扩展，提供本地HTTP服务器，支持文件保存后自动刷新浏览器。安装和使用步骤：

\begin{enumerate}[leftmargin=2em]
  \item 在VS Code中安装"Live Server"扩展
  \item 右键点击\texttt{web/index.html}，选择"Open with Live Server"
  \item 浏览器自动打开\texttt{http://127.0.0.1:5500/index.html}
  \item 修改代码并保存后，浏览器自动刷新
\end{enumerate}

\subsection{可选进阶：部署到公网}

如果你希望将CRM系统部署到公网，让别人也能访问，以下是两种零成本的方案：

\textbf{方案一：EdgeOne Pages}。EdgeOne Pages是腾讯云提供的全球CDN静态网站托管服务，操作步骤：
\begin{enumerate}[leftmargin=2em]
  \item 访问\url{https://console.cloud.tencent.com/edgeone/pages}，创建「直接上传」类型项目
  \item 获取API Token（\url{https://pages.edgeone.ai/document/api-token}）
  \item 使用CLI部署：\texttt{npx edgeone pages deploy ./web -n my-crm -t \$EDGEONE\_PAGES\_API\_TOKEN}
  \item 获得\texttt{xxx.edgeone.app}域名，国内访问速度快
\end{enumerate}

\textbf{方案二：Netlify}。Netlify支持拖拽部署，无需Git：
\begin{enumerate}[leftmargin=2em]
  \item 访问\url{https://app.netlify.com/drop}
  \item 将\texttt{web/}文件夹拖到页面上
  \item 自动部署，获得\texttt{xxx.netlify.app}域名
\end{enumerate}

\textbf{方案三：Cloudflare Pages}。Cloudflare Pages提供全球CDN和\texttt{wrangler} CLI直传部署：
\begin{enumerate}[leftmargin=2em]
  \item 安装CLI：\texttt{npm install -g wrangler}
  \item 登录账号：\texttt{wrangler login}（浏览器授权）
  \item 部署目录：\texttt{wrangler pages deploy ./web --project-name=my-crm --branch=main}
  \item 获得\texttt{xxx.pages.dev}域名，全球节点覆盖
\end{enumerate}

\begin{warningbox}[部署前检查清单]
部署到公网之前，请确认以下事项：
\begin{itemize}[leftmargin=2em]
  \item 所有功能在本地测试通过
  \item CDN资源链接使用HTTPS协议（\texttt{https://}而非\texttt{http://}）
  \item 无硬编码的localhost地址
  \item 模拟数据不包含真实个人信息（脱敏处理）
  \item 页面标题和描述已修改为你的项目信息
\end{itemize}
特别提醒：本CRM系统使用localStorage存储数据，这意味着数据存储在用户浏览器中，不同设备间的数据不共享。如果需要跨设备数据同步，必须搭建后端服务。\end{warningbox}

\begin{theorybox}[静态网站部署的技术原理]
EdgeOne Pages、Vercel和Netlify之所以能免费部署静态网站，是因为静态网站只需要一个文件服务器——用户的浏览器下载HTML/JS/CSS文件后在本地执行，服务器不需要做任何计算。这种架构被称为JAMstack（JavaScript + APIs + Markup），是现代Web开发的重要趋势。\par\n\n本课程的单文件HTML应用就是一个典型的静态网站——所有逻辑都在浏览器端执行，没有后端服务器的概念。如果未来加入Flask后端，就需要使用云服务器（如阿里云ECS、腾讯云CVM）或容器服务（如Railway、Fly.io）来部署后端代码，成本和技术复杂度都会增加。\end{theorybox}

\begin{practicebox}[完整测试流程]
\textbf{任务}：按照测试清单完成银行CRM的全面测试，并生成测试报告。

\textbf{提示词}：

\begin{lstlisting}
请帮我生成银行CRM系统的测试报告模板，包含：
1. 测试环境信息（浏览器版本/测试时间/数据量）
2. 测试结果汇总表（通过/失败/阻塞数量）
3. 详细测试记录（每条测试项的预期/实际/结论）
4. 发现的问题清单（严重程度/复现步骤/修复建议）
5. 测试结论与改进建议
保存到 reports/test_report.md
\end{lstlisting}

\textbf{操作流程}：
\begin{enumerate}[leftmargin=2em]
  \item 逐条执行功能测试清单中的测试项
  \item 记录每项的实际结果
  \item 发现bug后，用标准格式描述给AI修复
  \item AI修复后执行回归测试
  \item 全部测试通过后，生成最终测试报告
\end{enumerate}
\end{practicebox}

% ============================================================
\section{课程大作业要求}
% ============================================================

\subsection{学生交付清单}

完成银行CRM系统开发后，需提交以下交付物：

\begin{table}[H]
\centering
\begin{tabular}{clll}
\toprule
\textbf{\#} & \textbf{交付物} & \textbf{形式} & \textbf{说明} \\
\midrule
1 & 功能清单 & features.md & >= 5 项 \\
2 & 需求文档 & requirements.md & >= 4 条 R-编号 \\
3 & ER 图 & figures/crm\_er.mmd & mermaid格式 \\
4 & 设计文档 & design.md & 技术选型 + 架构图 \\
5 & 前端代码 & web/index.html & 单文件可运行 \\
6 & 测试报告 & reports/test\_report.md & 测试结果 \\
\bottomrule
\end{tabular}
\caption{学生交付清单}
\end{table}

\textbf{验收标准}：index.html 能在浏览器中打开，能完成"客户列表查看 + 新增客户 + 产品推荐"。

\subsection{答辩要求}

每位学生需进行\textbf{5分钟演示 + 3分钟问答}：

\begin{itemize}[leftmargin=2em]
  \item \textbf{演示内容}（5分钟）：
  \begin{enumerate}[leftmargin=2em]
    \item 系统功能演示（1分钟）：从打开到完成核心操作的全流程
    \item 技术亮点介绍（2分钟）：推荐算法、数据可视化等
    \item BMAD过程回顾（1分钟）：B/M/A/D/V各阶段的产出与心得
    \item 创新点展示（1分钟）：超出基本要求的功能
  \end{enumerate}
  \item \textbf{问答环节}（3分钟）：教师针对代码实现、设计决策、BMAD理解提问
\end{itemize}

\subsection{进阶加分项}

在完成基本要求的基础上，以下进阶功能可获得额外加分：

\begin{table}[H]
\centering
\begin{tabular}{clp{6cm}l}
\toprule
\textbf{编号} & \textbf{进阶功能} & \textbf{实现要求} & \textbf{加分} \\
\midrule
B-001 & 登录功能 & 用户名密码验证 + 登录页面 + 会话管理 & +5分 \\
B-002 & 数据导入/导出 & 支持CSV/Excel格式导入客户数据 + 导出功能 & +5分 \\
B-003 & 深色模式 & 一键切换明暗主题，所有页面适配 & +3分 \\
B-004 & 部署到公网 & 部署到EdgeOne Pages/Netlify，可公网访问 & +5分 \\
B-005 & 编辑/删除客户 & 完整CRUD + 确认弹窗 + 逻辑删除 & +3分 \\
B-006 & 移动端适配 & 手机端底部导航 + 单列布局 + 触摸友好 & +3分 \\
\bottomrule
\end{tabular}
\caption{进阶加分项}
\end{table}

\begin{notebox}[进阶加分上限]
进阶加分累计不超过15分（总分100分封顶）。鼓励你在核心功能稳定的基础上再尝试进阶功能，不要为了加分而忽视基本功能的完成度。
\end{notebox}

% ============================================================
\section{本章小结}
% ============================================================

本章以银行CRM系统为载体，完整实践了BMAD方法论的五个阶段，并将第9章的理论转化为可运行的系统：

\begin{enumerate}[leftmargin=2em]
  \item \textbf{需求分析}：通过用户画像和用户故事明确了系统的功能范围和优先级
  \item \textbf{界面设计}：基于Tailwind CSS实现了银行蓝色系的响应式界面
  \item \textbf{客户管理}：实现了完整的CRUD功能和localStorage持久化
  \item \textbf{产品推荐}：构建了基于三条业务规则的推荐引擎
  \item \textbf{数据可视化}：使用Chart.js实现了4种核心图表和KPI看板
  \item \textbf{测试部署}：完成了功能测试清单和测试报告生成
\end{enumerate}

通过本章的实践，你应该已经掌握了\textbf{从需求到交付}的完整开发流程，以及使用AI辅助各阶段工作的方法。这套方法论不仅适用于银行CRM，也可以迁移到其他Web应用的开发中。

% ============================================================
\section{验收清单}
% ============================================================

请逐项检查你的CRM系统：

\begin{table}[H]
\centering
\begin{tabular}{clll}
\toprule
\textbf{序号} & \textbf{检查项} & \textbf{验证方法} & \textbf{期望结果} \\
\midrule
1 & 页面可打开 & Live Server打开index.html & 显示客户列表 \\
2 & 客户列表正常 & 查看列表 & 显示>=5条客户 \\
3 & 搜索功能正常 & 输入姓名关键词 & 返回匹配结果 \\
4 & VIP筛选正常 & 选择VIP等级 & 列表按等级过滤 \\
5 & 新增客户正常 & 填写完整表单提交 & 列表新增一条记录 \\
6 & 表单验证生效 & 姓名为空提交 & 提示"请输入姓名" \\
7 & 客户详情正常 & 点击列表某行 & 显示完整信息 \\
8 & 产品推荐正常 & 选择客户 & 显示Top 3推荐 \\
9 & 推荐理由显示 & 查看推荐卡片 & 显示匹配规则 \\
10 & 图表渲染正常 & 点击报表分析 & 4个图表正常显示 \\
11 & KPI卡片正确 & 查看数值 & 与列表数据一致 \\
12 & 数据持久化 & 新增后刷新 & 数据不丢失 \\
13 & 交付物齐全 & 检查文件 & 6个文件完整 \\
\bottomrule
\end{tabular}
\caption{第10章验收清单}
\end{table}

% ============================================================
\section{思考题}
% ============================================================

\begin{exercisebox}[思考与讨论]
\begin{enumerate}[leftmargin=2em]
  \item 本系统的产品推荐基于三条业务规则，这在实际银行中是否足够？如果需要更精准的推荐，你会引入什么方法？
  \item localStorage存储的数据在清除浏览器缓存后会丢失。如果要求"数据永不丢失"，你会如何改造数据层？需要修改哪些代码？
  \item 本系统的推荐话术是预设的模板。如果要让AI根据客户的具体情况动态生成推荐话术，你会如何设计？需要注意什么问题？
  \item 如果你是一名银行客户经理，最希望CRM系统增加什么功能？请用用户故事格式描述。
  \item 将BMAD方法应用于本项目，你遇到了哪些困难？如果重新开始，你会如何改进工作流程？
\end{enumerate}
\end{exercisebox}
